%% parts/p02/first-order-normal-forms-supplement.tex --- retained additions
\section*{Editorial normal-form additions}
\addcontentsline{toc}{section}{Editorial normal-form additions}
The preceding normal-form definitions and proofs fix the formula language,
capture-avoiding substitution, and nonempty-domain semantics used here.
The quantifier-law table retrieves their side conditions.  The clause
construction below supplies the matrix conversion and the composed
expansion--reduct argument needed to pass from sentences to clause sets.

\begin{hbtheorem}{thm:p02-fol-quantifier-equivalences}{proved=annex:p2.fol.clausal.quantifier-laws}{Quantifier equivalence laws}
Every row of Table~\ref{tab:p02-quantifier-equivalences} is valid in the
nonempty-domain first-order semantics.  A row carrying
\(x\notin\FV(\chi)\) may be used only under that displayed side condition.
\end{hbtheorem}

\begingroup\small\setlength{\tabcolsep}{4pt}
\begin{longtable}{@{}p{.32in}p{2.55in}p{1.88in}@{}}
\caption{Quantifier equivalence laws with side conditions.}
\label{tab:p02-quantifier-equivalences}\\
\toprule
Row & Equivalence & Side condition \\ \midrule
\endfirsthead
\caption[]{Quantifier equivalence laws with side conditions (continued)}\\
\toprule Row & Equivalence & Side condition \\ \midrule
\endhead
\midrule\multicolumn{3}{r}{\emph{continued on next page}}\\
\endfoot
\endlastfoot
Q1 & \(\neg\forall x\,\varphi\Leftrightarrow
       \exists x\,\neg\varphi\) & none \\
Q2 & \(\neg\exists x\,\varphi\Leftrightarrow
       \forall x\,\neg\varphi\) & none \\
Q3 & \((Qx\,\varphi)\land\chi\Leftrightarrow
       Qx(\varphi\land\chi)\), \(Q\in\{\forall,\exists\}\) &
       \(x\notin\FV(\chi)\) \\
Q4 & \((Qx\,\varphi)\lor\chi\Leftrightarrow
       Qx(\varphi\lor\chi)\), \(Q\in\{\forall,\exists\}\) &
       \(x\notin\FV(\chi)\) \\
Q5 & \((\forall x\,\varphi)\to\chi\Leftrightarrow
       \exists x(\varphi\to\chi)\) & \(x\notin\FV(\chi)\) \\
Q6 & \((\exists x\,\varphi)\to\chi\Leftrightarrow
       \forall x(\varphi\to\chi)\) & \(x\notin\FV(\chi)\) \\
Q7 & \(\chi\to(\forall x\,\varphi)\Leftrightarrow
       \forall x(\chi\to\varphi)\) & \(x\notin\FV(\chi)\) \\
Q8 & \(\chi\to(\exists x\,\varphi)\Leftrightarrow
       \exists x(\chi\to\varphi)\) & \(x\notin\FV(\chi)\) \\
Q9 & \(\forall x(\varphi\land\psi)\Leftrightarrow
       (\forall x\varphi)\land(\forall x\psi)\) & none \\
Q10 & \(\exists x(\varphi\lor\psi)\Leftrightarrow
         (\exists x\varphi)\lor(\exists x\psi)\) & none \\
Q11 & \(\forall x\,\varphi\Leftrightarrow\varphi\) &
        \(x\notin\FV(\varphi)\) \\
Q12 & \(\exists x\,\varphi\Leftrightarrow\varphi\) &
        \(x\notin\FV(\varphi)\) \\
\bottomrule
\end{longtable}
\endgroup

The converses of the familiar one-way distribution laws omitted from the
table are generally invalid; for example, two existential conjuncts need not
share a witness.  The side conditions in Q3--Q8 are part of the law, not
typographical annotations.


\section{The clausal-form pipeline}
\label{sec:p02-fol-clausal-pipeline}

\begin{hbdefinition}{def:p02-clause-family-reading}{Clauses and their universal reading}
A clause \(C\) is a finite set of first-order literals.  At a structure
\(\mathcal N\) and assignment \(s\), it is true exactly when some
\(L\in C\) is true.  Thus the empty clause is false.  A finite clause
family \(S\) is true at \((\mathcal N,s)\) when every member is true;
the empty family is true.  A structure is a model of \(S\) when each
clause is true under every assignment to its free variables.  Equivalently,
it satisfies the universal closure of each clause's disjunction.
Duplicate literals and duplicate clauses do not affect these truth values.
The constants \(\bottom\) and \(\neg\bottom\) are treated as truth
units during conversion, not added to the first-order literal vocabulary.
\end{hbdefinition}

\begin{hbmethod}{def:p02-matrix-clause-construction}{Distributive clause construction}
For a quantifier-free NNF matrix \(\theta\), define a finite clause
family \(\operatorname{Cl}(\theta)\) by recursion:
\[
\begin{aligned}
 \operatorname{Cl}(L)&=\{\{L\}\} &&\text{for a literal }L,\\
 \operatorname{Cl}(\bottom)&=\{\varnothing\},&
 \operatorname{Cl}(\neg\bottom)&=\varnothing,\\
 \operatorname{Cl}(\alpha\land\beta)
   &=\operatorname{Cl}(\alpha)\cup\operatorname{Cl}(\beta),\\
 \operatorname{Cl}(\alpha\lor\beta)
   &=\{C\cup D:C\in\operatorname{Cl}(\alpha),\
                 D\in\operatorname{Cl}(\beta)\}.
\end{aligned}
\]
Each call is on an immediate subformula.  The construction preserves
literals as formula strings, including their terms and any equality
atoms.  It does not replace first-order atoms by new predicate symbols.
Repeated distribution can make the family exponentially large; no
linear-size bound is claimed for this conversion.
\end{hbmethod}

\begin{hblemma}{thm:p02-matrix-clause-equivalence}{proved=annex:p2.fol.clausal.matrix}{Pointwise correctness of matrix conversion}
For every nonempty-domain structure \(\mathcal N\), assignment \(s\),
and quantifier-free NNF matrix \(\theta\),
\[
 \mathcal N,s\entails\theta
 \quad\Longleftrightarrow\quad
 \text{every }C\in\operatorname{Cl}(\theta)
       \text{ contains an }L\text{ with }\mathcal N,s\entails L.
\]
The output is finite and introduces no variable or nonlogical symbol.
The equivalence includes the empty clause and empty family conventions.
\end{hblemma}

\begin{hbtheorem}{thm:p02-fol-clausal-form}{proved=annex:p2.fol.clausal.composition}{Clausal-form theorem}
For every first-order \(\mathcal L\)-sentence \(\sigma\), there is a finite set \(S\) of
universally read clauses in a finite expansion \(\mathcal L^*\) of its signature such that,
assuming AC, \(\sigma\) is satisfiable iff \(S\) is satisfiable.  The reverse
direction from an \(S\)-model to a model of \(\sigma\) uses only reducts and
does not require choice.  The transformation is the composition in
Table~\ref{tab:p02-fol-normal-form-pipeline}, following its CNF branch
(step 6 is a comparison, not an additional pipeline stage); only the Skolem step weakens
same-signature logical equivalence to expansion--reduct equisatisfiability.
\end{hbtheorem}

\begin{landscape}
\begingroup\scriptsize\setlength{\tabcolsep}{3pt}
\begin{longtable}{@{}p{.30in}p{1.05in}p{1.65in}p{1.90in}p{2.00in}@{}}
\caption{First-order normal-form comparison and clausal pipeline.}
\label{tab:p02-fol-normal-form-pipeline}\\
\toprule
Step & Output & Existence and scope & Preservation claim & Worst-case size and new data \\ \midrule
\endfirsthead
\caption[]{First-order normal-form comparison and clausal pipeline (continued)}\\
\toprule Step & Output & Existence and scope & Preservation claim & Worst-case size and new data \\ \midrule
\endhead
\midrule\multicolumn{5}{r}{\emph{continued on next page}}\\
\endfoot
\endlastfoot
1 & NNF & Exists for every formula in the same signature. & Formula
equivalence; exactly the same free and all-variable sets. & Linear in the
formula tree; no new symbols. \\
2 & Rectified NNF & Exists because the variable supply is countably infinite;
fresh binders are chosen outside a finite forbidden set. & Alpha-equivalence,
hence formula equivalence; same constructor tree. & Linear apart from the
encoding of fresh variable names; no new nonlogical symbols. \\
3 & Prenex NNF & Exists for every rectified NNF formula by the side-conditioned
laws in Table~\ref{tab:p02-quantifier-equivalences}. & Formula equivalence;
same free and all-variable sets and the same quantifier occurrences. & Linear
in the formula tree; only the prefix is rearranged. \\
4 & Universal Skolem sentence & Exists for every rectified prenex sentence in a
finite signature expansion; AC supplies witness functions in the forward
model-expansion direction. & Expansion--reduct equisatisfiability under AC,
not same-signature formula equivalence. & At most quadratic under literal term
printing because a new Skolem term may list all preceding universal variables;
new function symbols are introduced. \\
5 & Matrix in CNF & Exists for every quantifier-free matrix by propositional
distribution. & Formula equivalence in the expanded signature. & Direct
distribution can be exponential; Tseitin conversion is only equisatisfiable
and introduces symbols (see Part~2.1). \\
6 & Matrix in DNF & Exists for every quantifier-free matrix by dual
propositional distribution. & Formula equivalence in the expanded signature. &
Direct distribution can be exponential; DNF is a comparison form, not the
input used by resolution. \\
7 & Clause set & Exists from a universal sentence whose matrix is in CNF. & The
set and the universal closures of its clauses have exactly the same models as
that universal sentence. & Linear once the CNF matrix is present; erases the
written universal prefix and records a finite literal set per conjunct. \\
\bottomrule
\end{longtable}
\endgroup
\end{landscape}

\begin{figure}[htbp]
\centering
\begin{tikzpicture}[
  font=\small,
  node distance=8mm,
  box/.style={draw,rounded corners,align=center,inner sep=4pt,
              minimum width=1.55in},
  >=Stealth]
  \node[box] (in) {sentence\\\(\sigma\)};
  \node[box,below=of in] (nnf) {NNF};
  \node[box,below=of nnf] (rect) {rectified NNF};
  \node[box,below=of rect] (pren) {rectified prenex NNF};
  \node[box,below=of pren] (sk) {universal Skolem sentence};
  \node[box,below=of sk] (cnf) {universal sentence\\with CNF matrix};
  \node[box,below=of cnf] (clauses) {clause set \(S\)};
  \draw[->] (in) -- node[right=8mm,align=left,font=\footnotesize]
    {formula equivalence} (nnf);
  \draw[->] (nnf) -- node[right=8mm,align=left,font=\footnotesize]
    {alpha-equivalence} (rect);
  \draw[->] (rect) -- node[right=8mm,align=left,font=\footnotesize]
    {formula equivalence} (pren);
  \draw[->] (pren) -- node[right=8mm,align=left,font=\footnotesize]
    {equisatisfiable\\expanded signature} (sk);
  \draw[->] (sk) -- node[right=8mm,align=left,font=\footnotesize]
    {matrix CNF\\formula equivalence} (cnf);
  \draw[->] (cnf) -- node[right=8mm,align=left,font=\footnotesize]
    {universal-closure convention} (clauses);
\end{tikzpicture}
\caption{Normal-form pipeline.  The Skolem arrow records
expansion--reduct equisatisfiability rather than same-signature formula
equivalence; the last arrow relates a sentence to a clause set through the
universal closures of its members, not through formula-string equivalence.}
\label{fig:p02-fol-normal-form-pipeline}
\end{figure}

The complete composition proof is
Annex~\ref{annex:p2.fol.clausal.composition}; its matrix step is
Annex~\ref{annex:p2.fol.clausal.matrix}.  Universal reading is essential:
erasing quantifier symbols does not turn their bound variables into freely
chosen witnesses.  Resolution operations on these clause objects are
defined in the next chapter.  The table does not identify a finite clause
set with a single formula string.

\begin{hbexample}{def:p02-clausal-composition-example}{An alternating prefix and two clauses}
For
\[
 \forall x\exists y\,\bigl(R(x,y)\lor(P(x)\land Q(y))\bigr),
\]
the existing prefix is rectified and the matrix is in NNF.
Skolemization introduces one fresh unary \(f\).  Distribution gives
the clause family
\[
 \bigl\{\{R(x,f(x)),P(x)\},\
         \{R(x,f(x)),Q(f(x))\}\bigr\}.
\]
Both clauses are universally read in \(x\).  A model of the original
sentence expands by choosing a witness \(f(a)\) for every domain element
\(a\); a model of the clauses supplies that witness directly when
recovering the existential sentence in the reduct.
If the input is instead \(\neg\bottom\), the output is the empty
family; if it is \(\bottom\), the output is a family containing the
empty clause.  These two empty objects have opposite satisfiability
statuses.  Applying the sentence theorem to an open formula requires
first specifying its closure: \(P(x)\) read under a fixed assignment,
\(\exists xP(x)\), and \(\forall xP(x)\) are different inputs.
\end{hbexample}
For the quantifier equivalences and the clause-conversion pipeline,
respectively, see \cite{Enderton2001,BaazEglyLeitsch2001}.
